\documentclass[10pt,a4paper]{amsart}
\usepackage{amsmath,amssymb,amsthm,graphicx}
\usepackage{verbatim}
\newenvironment{reference}{\comment}{\endcomment}
%\newenvironment{reference}{}{}
\newenvironment{slogan}{\comment}{\endcomment}
\newenvironment{history}{\comment}{\endcomment}

% For commutative diagrams we use Xy-pic
\usepackage[all]{xy}

% We use 2cell for 2-commutative diagrams.
\xyoption{2cell}
% Theorem environments.
%
\theoremstyle{plain}
\newtheorem{theorem}[subsection]{Theorem}
\newtheorem{proposition}[subsection]{Proposition}
\newtheorem{lemma}[subsection]{Lemma}

\theoremstyle{definition}
\newtheorem{definition}[subsection]{Definition}
\newtheorem{example}[subsection]{Example}
\newtheorem{exercise}[subsection]{Exercise}
\newtheorem{situation}[subsection]{Situation}

\theoremstyle{remark}
\newtheorem{remark}[subsection]{Remark}
\newtheorem{remarks}[subsection]{Remarks}

\numberwithin{equation}{subsection}

% Macros
%
\def\lim{\mathop{\mathrm{lim}}\nolimits}
\def\colim{\mathop{\mathrm{colim}}\nolimits}
\def\Spec{\mathop{\mathrm{Spec}}}
\def\Hom{\mathop{\mathrm{Hom}}\nolimits}
\def\Ext{\mathop{\mathrm{Ext}}\nolimits}
\def\SheafHom{\mathop{\mathcal{H}\!\mathit{om}}\nolimits}
\def\SheafExt{\mathop{\mathcal{E}\!\mathit{xt}}\nolimits}
\def\Sch{\mathit{Sch}}
\def\Mor{\mathop{\mathrm{Mor}}\nolimits}
\def\Ob{\mathop{\mathrm{Ob}}\nolimits}
\def\Sh{\mathop{\mathit{Sh}}\nolimits}
\def\NL{\mathop{N\!L}\nolimits}
\def\CH{\mathop{\mathrm{CH}}\nolimits}
\def\proetale{{pro\text{-}\acute{e}tale}}
\def\etale{{\acute{e}tale}}
\def\QCoh{\mathit{QCoh}}
\def\Ker{\mathop{\mathrm{Ker}}}
\def\Im{\mathop{\mathrm{Im}}}
\def\Coker{\mathop{\mathrm{Coker}}}
\def\Coim{\mathop{\mathrm{Coim}}}

% Boxtimes
%
\DeclareMathSymbol{\boxtimes}{\mathbin}{AMSa}{"02}

%
% Macros for moduli stacks/spaces
%
\def\QCohstack{\mathcal{QC}\!\mathit{oh}}
\def\Cohstack{\mathcal{C}\!\mathit{oh}}
\def\Spacesstack{\mathcal{S}\!\mathit{paces}}
\def\Quotfunctor{\mathrm{Quot}}
\def\Hilbfunctor{\mathrm{Hilb}}
\def\Curvesstack{\mathcal{C}\!\mathit{urves}}
\def\Polarizedstack{\mathcal{P}\!\mathit{olarized}}
\def\Complexesstack{\mathcal{C}\!\mathit{omplexes}}
% \Pic is the operator that assigns to X its picard group, usage \Pic(X)
% \Picardstack_{X/B} denotes the Picard stack of X over B
% \Picardfunctor_{X/B} denotes the Picard functor of X over B
\def\Pic{\mathop{\mathrm{Pic}}\nolimits}
\def\Picardstack{\mathcal{P}\!\mathit{ic}}
\def\Picardfunctor{\mathrm{Pic}}
\def\Deformationcategory{\mathcal{D}\!\mathit{ef}}
\usepackage{cleveref}
\crefname{originalreference}{original reference}{original references}
\Crefname{originalreference}{Original reference}{Original references}
\makeatletter
\newlabel{situation-shriek}{{0F42}{source}{}{sourcepointer-1}{}}
\newlabel{situation-shriek@cref}{{[originalreference][0][]0F42}{[1][0][]source}{}{}{}}
\makeatother
\crefname{originalequation}{original equation}{original equations}
\Crefname{originalequation}{Original equation}{Original equations}
\makeatletter
\newlabel{equation-base-change-map}{{0AA6}{source}{}{sourcepointer-2}{}}
\newlabel{equation-base-change-map@cref}{{[originalequation][0][]0AA6}{[1][0][]source}{}{}{}}
\makeatother
\crefname{originallemma}{original lemma}{original lemmas}
\Crefname{originallemma}{Original lemma}{Original lemmas}
\makeatletter
\newlabel{lemma-twisted-inverse-image}{{0A9E}{source}{}{sourcepointer-3}{}}
\newlabel{lemma-twisted-inverse-image@cref}{{[originallemma][0][]0A9E}{[1][0][]source}{}{}{}}
\makeatother
\crefname{originallemma}{original lemma}{original lemmas}
\Crefname{originallemma}{Original lemma}{Original lemmas}
\makeatletter
\newlabel{lemma-proper-noetherian}{{0A9P}{source}{}{sourcepointer-4}{}}
\newlabel{lemma-proper-noetherian@cref}{{[originallemma][0][]0A9P}{[1][0][]source}{}{}{}}
\makeatother
\crefname{originalreference}{original reference}{original references}
\Crefname{originalreference}{Original reference}{Original references}
\makeatletter
\newlabel{remark-check-over-affines}{{0E9S}{source}{}{sourcepointer-5}{}}
\newlabel{remark-check-over-affines@cref}{{[originalreference][0][]0E9S}{[1][0][]source}{}{}{}}
\makeatother
\crefname{originallemma}{original lemma}{original lemmas}
\Crefname{originallemma}{Original lemma}{Original lemmas}
\makeatletter
\newlabel{lemma-more-base-change}{{0AA8}{source}{}{sourcepointer-6}{}}
\newlabel{lemma-more-base-change@cref}{{[originallemma][0][]0AA8}{[1][0][]source}{}{}{}}
\makeatother
\title{Base change for the upper shriek along a flat morphism}
\author{The Stacks Project Authors}
\date{stacks-project:duality.tex:local-snapshot-20261012; GNU FDL 1.2 or later}
\begin{document}
\maketitle
\noindent\textit{Editorial note.} All delivered mathematical text is the author's own, copied from the pinned Stacks Project chapter file duality.tex; no statement or proof step was rewritten, reordered or removed.\par\medskip
\noindent\textit{Editorial note.} Preamble provenance: the pinned chapter file opens with an input of the shared preamble and carries no preamble text of its own, so the project's single Stacks-wide preamble staged at sources/stacks-duality/preamble.tex is excerpted verbatim under the recipe fixed by LANE\_BRIEF 22.1, as three spans, lines 8-13 (verbatim package and the reference/slogan/history environments), lines 14-18 (Xy-pic setup) and lines 154-217 (the native theorem declarations and mathematical macros). Nothing was added, removed or redefined; the wrapper class and the amsmath/amssymb/amsthm/graphicx base packages are supplied by the assembler.\par\medskip
\noindent\textit{Editorial note.} The counted statement and proof contain no cross reference, no citation and no included file, so the delivered body is self contained.\par\medskip
\noindent\textit{Editorial note.} Canonical location of the counted proposition: Stacks tag 0E9U (sources/stacks-duality/pinned-original-source.tex line 4997).\par\medskip
\noindent\textit{Allowed standard prerequisites (editorial).}\begin{itemize}
\item The upper shriek through a compactification and the base change map for the right adjoints. Same chapter, Lemma labelled lemma-twisted-inverse-image and equation labelled equation-base-change-map, cited at proof lines 5015-5028. Define the map whose bijectivity is asserted.
\item Affine locality of the right adjoint and the reduction to affine targets. Same chapter, Lemma labelled lemma-proper-noetherian and Remark labelled remark-check-over-affines, cited at proof lines 5036-5040. Reduce the isomorphism to affine Y and Y'.
\item The base change statement in the affine flat case. Same chapter, Lemma labelled lemma-more-base-change, cited at proof lines 5040-5041. Settles the reduced case.
\end{itemize}\medskip
\ifdefined\hypertarget\hypertarget{sourcepointer-1}{}\fi
\noindent\textit{Source reference.} Original reference 0F42: sources/stacks-duality/pinned-original-source.tex, source lines 3953--3953. Same-chapter pointer into the author environment situation at source line 3953; invoked at proof line None of the counted proof. It is retained only as a navigation pointer; every load-bearing step is inside the delivered span.\par\medskip
\ifdefined\hypertarget\hypertarget{sourcepointer-2}{}\fi
\noindent\textit{Source reference.} Original equation 0AA6: sources/stacks-duality/pinned-original-source.tex, source lines 1192--1192. Same-chapter pointer into the author environment equation at source line 1192; invoked at proof line 5024 of the counted proof. It is retained only as a navigation pointer; every load-bearing step is inside the delivered span.\par\medskip
\ifdefined\hypertarget\hypertarget{sourcepointer-3}{}\fi
\noindent\textit{Source reference.} Original lemma 0A9E: sources/stacks-duality/pinned-original-source.tex, source lines 653--653. Same-chapter pointer into the author environment lemma at source line 653; invoked at proof line 5026 of the counted proof. It is retained only as a navigation pointer; every load-bearing step is inside the delivered span.\par\medskip
\ifdefined\hypertarget\hypertarget{sourcepointer-4}{}\fi
\noindent\textit{Source reference.} Original lemma 0A9P: sources/stacks-duality/pinned-original-source.tex, source lines 1106--1106. Same-chapter pointer into the author environment lemma at source line 1106; invoked at proof line 5036 of the counted proof. It is retained only as a navigation pointer; every load-bearing step is inside the delivered span.\par\medskip
\ifdefined\hypertarget\hypertarget{sourcepointer-5}{}\fi
\noindent\textit{Source reference.} Original reference 0E9S: sources/stacks-duality/pinned-original-source.tex, source lines 1585--1585. Same-chapter pointer into the author environment remark at source line 1585; invoked at proof line 5038 of the counted proof. It is retained only as a navigation pointer; every load-bearing step is inside the delivered span.\par\medskip
\ifdefined\hypertarget\hypertarget{sourcepointer-6}{}\fi
\noindent\textit{Source reference.} Original lemma 0AA8: sources/stacks-duality/pinned-original-source.tex, source lines 1636--1636. Same-chapter pointer into the author environment lemma at source line 1636; invoked at proof line 5040 of the counted proof. It is retained only as a navigation pointer; every load-bearing step is inside the delivered span.\par\medskip
\begin{lemma}
\label{lemma-base-change-shriek-flat}
In Situation \ref{situation-shriek} let
$$
\xymatrix{
X' \ar[r]_{g'} \ar[d]_{f'} & X \ar[d]^f \\
Y' \ar[r]^g & Y
}
$$
be a cartesian diagram of $\textit{FTS}_S$ with $g$ flat.
Then there is an isomorphism
$L(g')^* \circ f^! \to (f')^! \circ Lg^*$ on
$D_\QCoh^+(\mathcal{O}_Y)$.
\end{lemma}

\begin{proof}
Namely, because $g$ is flat, for every choice of
compactification $j : X \to \overline{X}$ of $X$ over $Y$
the scheme $\overline{X}$ is Tor independent of $Y'$.
Denote $j' : X' \to \overline{X}'$ the
base change of $j$ and $\overline{g}' : \overline{X}' \to \overline{X}$
the projection. We define the base change map as the composition
$$
L(g')^* \circ f^! = L(g')^* \circ j^* \circ a =
(j')^* \circ L(\overline{g}')^* \circ a \longrightarrow
(j')^* \circ a' \circ Lg^* = (f')^! \circ Lg^*
$$
where the middle arrow is the base change map
(\ref{equation-base-change-map})
and $a$ and $a'$ are the right adjoints to pushforward
of Lemma \ref{lemma-twisted-inverse-image}
for $\overline{X} \to Y$ and $\overline{X}' \to Y'$.
This construction is independent of the choice of
compactification (we will formulate a precise lemma
and prove it, if we ever need this result).

\medskip\noindent
To finish the proof it suffices to show that the base change
map $L(g')^* \circ a \to a' \circ Lg^*$ is an isomorphism
on $D_\QCoh^+(\mathcal{O}_Y)$.
By Lemma \ref{lemma-proper-noetherian} formation of $a$ and $a'$
commutes with restriction to affine opens of $Y$ and $Y'$.
Thus by Remark \ref{remark-check-over-affines}
we may assume that $Y$ and $Y'$ are affine.
Thus the result by Lemma \ref{lemma-more-base-change}.
\end{proof}

\end{document}
